\documentclass{axiobac}

\annee{1991}
\mois{Juin}
\session{principale}
\epreuve{Mathématiques}
\sectionbac{Math.-Sc. (Sp. : Mathématiques) \& Mathématiques Technique}
\sections{mathematiques, sciences-techniques}
\duree{4 heures}
\coefficient{5}
\ministere{Ministère de l'Éducation et des Sciences}

\begin{document}
\entete

%----------------------------------------------------------------------------------------
%   EXERCICE 1
%----------------------------------------------------------------------------------------
\begin{exercice}{6}
\titre{Triangle équilatéral, cercles tangents et rotations}
\chapitres{Déplacements - Antidéplacements}

On considère, dans un plan orienté, un triangle équilatéral $ABC$ tel que $\left(\widehat{\overrightarrow{AB}, \overrightarrow{AC}}\right) \equiv \dfrac{\pi}{3} \ [2\pi]$ et le cercle $(\mathcal{C})$ circonscrit à ce triangle. On construit les cercles $(\mathcal{C}')$ et $(\mathcal{C}'')$ de même rayon, de centres respectifs $B$ et $C$ et tangents extérieurement l'un à l'autre.

\begin{enumerate}
    \item Soit $M$ un point du cercle $(\mathcal{C}')$ et $M'$ le point du cercle $(\mathcal{C}'')$ tel que $\left(\widehat{\overrightarrow{BM}, \overrightarrow{CM'}}\right) \equiv \dfrac{\pi}{3} \ [2\pi]$. \\
    Montrer que la médiatrice du segment $[MM']$ passe par un point fixe que l'on précisera.
    \item Le point $M$ se projette orthogonalement en $H$ sur $(AM')$ ; déterminer l'ensemble des points $H$ lorsque $M$ décrit $(\mathcal{C}')$. \\
    Construire cet ensemble.
    \item Soit $D$ le point du cercle $(\mathcal{C})$ diamétralement opposé à $A$ et $M''$ l'image du point $M'$ par la rotation de centre $D$ et d'angle dont une mesure est $\dfrac{2\pi}{3}$. \\
    Montrer que $M$ et $M''$ sont deux points diamétralement opposés du cercle $(\mathcal{C}')$.
\end{enumerate}

\end{exercice}

%----------------------------------------------------------------------------------------
%   EXERCICE 2
%----------------------------------------------------------------------------------------
\begin{exercice}{4}
\titre{Jetons blancs et noirs : tirages avec et sans remise}
\chapitres{Probabilités}

Un sac contient sept jetons indiscernables au toucher et répartis comme suit :
\begin{itemize}
    \item quatre jetons blancs numérotés $1$, $2$, $2$, $2$ ;
    \item trois jetons noirs numérotés $1$, $1$, $2$.
\end{itemize}

\begin{enumerate}
    \item On tire successivement et avec remise trois jetons du sac.
    \begin{enumerate}
        \item Calculer la probabilité d'avoir trois jetons blancs.
        \item Calculer la probabilité pour que, parmi les trois jetons tirés, il y en ait deux seulement qui portent le numéro $1$.
    \end{enumerate}
    \item On remet tous les jetons dans le sac et on tire de nouveau et successivement trois jetons de la manière suivante :
    \begin{itemize}
        \item si le jeton tiré porte le numéro $2$, il est remis dans le sac ;
        \item si le jeton tiré porte le numéro $1$, il n'est pas remis dans le sac.
    \end{itemize}
    \begin{enumerate}
        \item Calculer la probabilité d'avoir trois jetons blancs.
        \item Calculer la probabilité pour que deux seulement des trois jetons tirés portent le numéro $1$.
    \end{enumerate}
\end{enumerate}

\textit{N.B. : Toutes les probabilités seront calculées à $0{,}01$ près.}

\end{exercice}

%----------------------------------------------------------------------------------------
%   PROBLÈME
%----------------------------------------------------------------------------------------
\begin{probleme}{10}

Soit $f$ la fonction numérique à variable réelle définie par
\[
f(x) = \operatorname{Log}\left( \frac{x + 2 + \sqrt{x^2 + 4x}}{2} \right) .
\]

\partie{A}
\titre{Une fonction logarithme et sa réciproque $e^x + e^{-x} - 2$}
\chapitres{Fonction logarithme népérien, Fonctions réciproques, Fonctions exponentielles}

\begin{enumerate}
    \item
    \begin{enumerate}
        \item Montrer que le domaine de définition de $f$ est $\mathbb{R}_+$.
        \item Étudier les variations de $f$ sur $\mathbb{R}_+$.
        \item Tracer sa courbe représentative $(C)$ dans le plan rapporté à un repère orthonormé $(O, \vec{i}, \vec{j})$.
    \end{enumerate}
    \item
    \begin{enumerate}
        \item Soit $y$ un réel strictement positif. Résoudre dans $\mathbb{R}_+$ l'équation en $x$ : $e^x + e^{-x} - 2 = y$.
        \item Montrer que $f$ admet sur $\mathbb{R}_+$ une fonction réciproque $g$, et déduire de ce qui précède l'expression de $g(x)$.
    \end{enumerate}
\end{enumerate}

\partie{B}
\titre{Intégrale d'une fonction réciproque et aire}
\chapitres{Primitives, Intégrales, Fonctions réciproques}

\begin{enumerate}
    \item Soit $h$ une fonction numérique à variable réelle, dérivable et strictement monotone sur un intervalle $I$.
    \begin{enumerate}
        \item Montrer que $h^{-1}$ admet des primitives sur $h(I)$.
        \item Soit $H$ une primitive de $h^{-1}$ sur $h(I)$. Démontrer que $H \circ h$ est une primitive, sur $I$, de la fonction $x \longmapsto x\,h'(x)$.
        \item Soit $(\alpha, \beta)$ un couple de $I^2$. Déduire de ce qui précède que :
        \[
        \int_{h(\alpha)}^{h(\beta)} h^{-1}(t)\,dt = \int_\alpha^\beta t\,h'(t)\,dt .
        \]
    \end{enumerate}
    \item
    \begin{enumerate}
        \item Calculer $f(0)$ et $f\left(\dfrac{(e - 1)^2}{e}\right)$.
        \item En déduire l'aire de la partie du plan limitée par la courbe $(C)$ et les droites d'équations $x = \dfrac{(e - 1)^2}{e}$ et $y = 0$.
    \end{enumerate}
\end{enumerate}

\partie{C}
\titre{Encadrement et suite récurrente $u_{n+1} = f(2u_n)$}
\chapitres{Suites réelles, Fonction logarithme népérien}

\begin{enumerate}
    \item
    \begin{enumerate}
        \item Montrer que pour tout réel $x$ positif, on a : $\operatorname{Log}(1 + x) \leqslant f(x) \leqslant \operatorname{Log}(2 + x)$.
        \item Soit $\theta$ la fonction définie sur $\mathbb{R}_+$ par $\theta(x) = f(x) - \dfrac{1}{2}x$. \\
        Montrer que $\theta$ est strictement décroissante sur $[1, +\infty[$ et en déduire que l'équation $f(x) = \dfrac{1}{2}x$ admet une solution unique $\gamma$ dans $]1, 4[$.
    \end{enumerate}
    \item On considère la suite $(u_n)_{n \in \mathbb{N}}$ définie par :
    \[
    \begin{cases}
        u_0 = 1 \\
        u_{n+1} = f(2u_n) \quad \text{pour tout } n \in \mathbb{N}
    \end{cases}
    \]
    \begin{enumerate}
        \item Montrer que, pour tout $n \in \mathbb{N}$, on a : $u_n \in [1, 2]$.
        \item Montrer que, pour tout $n \in \mathbb{N}$, on a : $\left|u_{n+1} - \dfrac{1}{2}\gamma\right| \leqslant \dfrac{2}{\sqrt{5}} \left|u_n - \dfrac{1}{2}\gamma\right|$.
        \item En déduire la limite quand $n \to +\infty$ de la suite $(u_n)_{n \in \mathbb{N}}$.
    \end{enumerate}
\end{enumerate}

\end{probleme}

\end{document}
